\section{Appendix D. Negative Results}

\subsection{D.1 Failed Attempts}

\subsubsection{D.1.1 Very Low Rank (r=8)}

\textbf{Result.} {{R8\_ACC}}\%. DCT truncation loses too much high-frequency info.

\subsubsection{D.1.2 Learned Basis}

\textbf{Goal.} Replace DCT with learned projection.

\textbf{Result.} No improvement, more parameters. DCT is already optimal for 
smooth attention.

\subsubsection{D.1.3 Applying Galerkin to Decoders}

\textbf{Result.} Causality broken. Output incoherent.

\subsubsection{D.1.4 No Energy Scaling}

\textbf{Result.} {{NO\_SCALE\_ACC}}\%. Output magnitudes drift.

\subsubsection{D.1.5 Very High Rank (r=256)}

\textbf{Result.} {{R256\_ACC}}\%. No accuracy benefit, slower.

\subsubsection{D.1.6 No QR Orthonormalization}

\textbf{Result.} Accumulating floating-point drift; slight accuracy drop.

\subsection{D.2 Hyperparameter Failures}

\begin{table}[H]
\centering
\begin{tabular}{ll}
\toprule
\textbf{LR} & \textbf{D acc} \\
\midrule
5e-5 & {{D\_LR5e5}} \\
1e-4 & {{D\_LR1e4}} \\
2e-4 & {{D\_ACC}} \\
5e-4 & {{D\_LR5e4}} \\
\bottomrule
\end{tabular}
\end{table}

\subsection{D.3 Multi-Seed}

\begin{table}[H]
\centering
\begin{tabular}{llll}
\toprule
\textbf{Seed} & \textbf{D acc} & \textbf{R acc} & \textbf{G acc} \\
\midrule
42 & {{SEED42\_D}} & {{SEED42\_R}} & {{SEED42\_G}} \\
123 & {{SEED123\_D}} & {{SEED123\_R}} & {{SEED123\_G}} \\
456 & {{SEED456\_D}} & {{SEED456\_R}} & {{SEED456\_G}} \\
Mean & {{MEAN\_D}} & {{MEAN\_R}} & {{MEAN\_G}} \\
Std & {{STD\_D}} & {{STD\_R}} & {{STD\_G}} \\
\bottomrule
\end{tabular}
\end{table}

\subsection{D.4 Lessons Learned}

1. Encoder-only benefit is real.
2. Decoder fallback necessary.
3. DCT beats learned basis.
4. Rank 64 is a sweet spot.
5. Energy scaling is essential.

\subsection{D.5 Unresolved Questions}

1. \textbf{Causal Galerkin} -- is it possible?
2. \textbf{Learned basis} -- could it help for specific tasks?
3. \textbf{Very long context} ($N > 4096$) -- does DCT scaling hold?